% Generated draft: requires mathematical and source review.
\documentclass[11pt]{article}
\usepackage[margin=25mm]{geometry}
\usepackage{fontspec}
\setmainfont{Latin Modern Roman}
\usepackage{amsmath,amssymb,mathtools,unicode-math}
\setmathfont{Latin Modern Math}
\usepackage{enumitem,needspace,xurl,hyperref}
\setlength{\parindent}{0pt}
\setlength{\parskip}{5pt}
\setlength{\emergencystretch}{3em}
\newcommand{\N}{\mathbb{N}}
\newcommand{\Z}{\mathbb{Z}}
\newcommand{\Q}{\mathbb{Q}}
\newcommand{\R}{\mathbb{R}}
\newcommand{\C}{\mathbb{C}}
\newcommand{\F}{\mathbb{F}}
\newcommand{\A}{\mathbb{A}}
\newcommand{\PP}{\mathbb{P}}
\newcommand{\E}{\mathbb{E}}
\newcommand{\eps}{\varepsilon}
\newcommand{\dd}{\,\mathrm{d}}
\begin{document}
\section*{UnsolvedMath ID: 2700005}
\textbf{Source problem code:} catalogue entry
\subsection{Definitions and mathematical statement}
Denis Serre's problem is to prove existence of a regular shock reflection for compressible flow against a wedge without imposing irrotationality. In planar terminology, irrotationality is the condition $\partial_1u_2-\partial_2u_1=0$; removing this assumption allows vorticity and does not require it to be nonzero everywhere.

A standard full-flow context is the compressible Euler system
\[
 \partial_t\rho+\operatorname{div}(\rho u)=0,
 \qquad
 \partial_t(\rho u)+\operatorname{div}(\rho u\otimes u)+\nabla p=0,
\]
\[
 \partial_t E+\operatorname{div}((E+p)u)=0,
 \qquad E=\rho e+\tfrac12\rho|u|^2,
\]
with a specified equation of state and appropriate shock jump/entropy conditions. At an impermeable wedge, the usual slip condition is $u\cdot n=0$. Regular reflection refers to the configuration with an incident and a reflected shock meeting at the wedge, rather than a Mach-reflection configuration with an additional stem/triple-point pattern. These equations and descriptions supply context, not omitted source data.

The record does not give the precise governing model, equation of state, incident states, wedge angle or admissible angle regime, initial/boundary conditions, steady versus unsteady/self-similar formulation, or weak-solution/entropy and regularity requirements. It therefore does not determine a universal existence assertion for arbitrary wedge data. The target is existence in Serre's intended regular-reflection regime for a model that permits rotational flow; a potential-flow reduction is not a substitute for this task.
\subsection{Short English statement}
Establish regular reflection of compressible flow at a wedge using the full rotational-flow setting rather than assuming a velocity potential. The supplied assessment reports rigorous potential-flow work by Chen--Feldman and unspecified 2024--25 extensions with vorticity/non-potential flow, but says existence in the full exact Serre formulation remains open. The missing formulation and incomplete references prevent a stronger precise theorem or independent status conclusion here.
\subsection{Sources}
\begin{itemize}
\item UnsolvedMath permanent ID 2700005, source code catalogue entry, ``Five Open Problems --- Regular reflection without irrotationality.'' Registry curation and metadata: Ulam AI and UnsolvedMath Contributors, CC BY 4.0; proposer Denis Serre, 2012; updated 2026-08-06.
\item Serre, source list \emph{Problems} (2012), Open Problem 5: \url{https://perso.ens-lyon.fr/denis.serre/DPF/Ouverts.pdf}. Supplied PDF-text extraction accessed 2026-07-29. The metadata reports inspection of the corrected author PDF and an open journal article, presenting five open problems; no journal citation is supplied. Current-status release review and rights review remain pending.
\item Chen--Feldman and recent 2024--25 work are cited in the supplied partial-progress summary without complete titles, theorem hypotheses or locators. No authors for the unspecified later work or missing references are inferred.
\end{itemize}
\end{document}
